Käesoleva bakalaureusetöö eesmärk on uurida, kuidas luua ja hinnata eestikeelset äratussõna tuvastust piiratud ressursiga nutikodu mikrokontrolleril. Töö keskmes on projekt \enquote{Kratt}, mille siht on fraasi \enquote{Kuule Kratt} lokaalne tuvastamine ESP32-S3 klassi seadmel ilma pilveteenust kasutamata.

Töös rekonstrueeriti \texttt{microWakeWord}il põhinev treeningu- ja hindamistoru ning valideeriti see avaliku \texttt{Speech Commands} kontrollkatsega. Esialgne klipitasemel hindamine osutus andmelekke tõttu eksitavaks, mistõttu koostati sõltumatud kõrvalejäetud testikomplektid, lisati treening- ja testandmete mittekattuvuse kontroll ning võeti keskseks mõõdikuks voogedastusrežiimi FAPH ehk valevallandumiste arv tunnis.

Tulemused näitavad, et ükski praegune mudel ega mudelite kombinatsioon ei täida korraga kõiki sihte: madalat FAPH-i, kõrget päriskõnelejate tuvastamismäära ja täpset fraasieristust. Töö peamine panus on seetõttu mitte tootmisküps mudel, vaid reprodutseeritav hindamisprotokoll ja empiiriline kirjeldus kompromissidest, mis tekivad väikese ressursiga keele lokaalse äratussõna arendamisel.

Lõputöö on kirjutatud \langEst~keeles ning sisaldab teksti \calculatepages leheküljel, 
\total{totalchapters} peatükki\ifthenelse{\equal{\totvalue{figure}}{0}}{}{% If no figures, do nothing
, \total{figure} joonist}%
\ifthenelse{\equal{\totvalue{table}}{0}}{}{% If no tables, do nothing
, \total{table} tabelit}%
.
